\lysh{35774}{2}
\begin{alist}
    \item Η ζητούμενη μέση τιμή της βαθμολογίας του μαθητή, για τα πέντε τεστ είναι:
    \[
        \bar{x} = \frac{\sum_{i=1}^{\nu} t_i}{\nu} = \frac{12+14+8+20+16}{5} = \frac{70}{5} = 14.
    \]

    \item
    \begin{rlist}
        \item Λόγω του ερωτήματος (α) είναι $\bar{x}=14$, οπότε ο πίνακας συμπληρωμένος είναι:
        \begin{center}
        \begin{mytblr}{}
            {\shortstack{Βαθμός \\ $x_i$}} & $x_i - \bar{x}$ & $(x_i - \bar{x})^2$ \\
            
            12 & -2 & 4 \\
            14 & 0 & 0 \\
            8 & -6 & 36 \\
            20 & 6 & 36 \\
            16 & 2 & 4 \\
            
            ΣΥΝΟΛΟ & \SetCell{bg=gray!25} {} & 80 \\
        \end{mytblr}
        \end{center}
        
        \item Λόγω του ερωτήματος (β, i) είναι:
        \[
            s^2 = \frac{\sum_{i=1}^{\nu} (x_i - \bar{x})^2}{\nu} = \frac{80}{5} = 16
        \]
        η διακύμανση της βαθμολογίας και
        \[
            s = \sqrt{s^2} = \sqrt{16} = 4
        \]
        η τυπική απόκλιση.
    \end{rlist}
\end{alist}
\vspace{6mm}
\noindent\rule{\textwidth}{0.4pt}
\vspace{6mm}

